\documentclass[UTF8,zihao=-4]{ctexart}
\usepackage[a4paper,margin=2.2cm]{geometry}
\usepackage{array}
\usepackage{longtable}
\usepackage{hyperref}
\usepackage{verbatim}
\hypersetup{hidelinks}
\setlength{\parindent}{0pt}
\setlength{\parskip}{0.55em}
\setcounter{secnumdepth}{0}
\renewcommand{\arraystretch}{1.25}
\emergencystretch=4em
\sloppy
\title{01 软件工程基础与项目管理复习讲义}
\author{}
\date{}
\begin{document}
\maketitle
\tableofcontents
\newpage


\begin{quote}
范围：\texttt{01-01-软件工程基础.pdf}、\texttt{01-02-项目管理基础.pdf}；结合 Exam/Review 中软件工程定义、SQA、SCM、变更控制和风险管理内容。
结构：先整理 PPT 主线，再放对应考试模板和快速背诵清单。
\end{quote}

\bigskip\hrule\bigskip

\section{1. 软件工程基础与项目管理}

\subsection{1.1 软件是什么}

软件不是单纯的程序。课程中的核心表达是：

\textbf{软件 = 程序 + 数据 + 文档 + 知识}

IEEE 610.12-1990 对软件的定义强调三类内容：

\begin{itemize}
\item 程序
\item 规程或过程
\item 相关文档与数据
\end{itemize}

答题时要写出“软件不只等于代码”。代码只是可执行部分，文档、数据、过程和领域知识共同决定软件能否被理解、运行、维护和演化。

\subsection{1.2 软件的典型特征}

常见概括：

\begin{itemize}
\item 无形：软件没有物理形态，进度和质量不易直接观察。
\item 复杂：状态、分支、交互和依赖数量高。
\item 可复制：复制成本低，但开发和维护成本高。
\item 易修改：修改入口多，但影响传播复杂。
\item 退化：软件不会物理磨损，但会因环境变化、需求变化和补丁叠加而退化。
\item 依赖：依赖硬件、操作系统、网络、数据库、第三方库和组织流程。
\item 智力密集：核心产出来自设计、抽象、分析和协作。
\end{itemize}

考试答“软件危机”时，要把这些特征与项目失败联系起来：规模变大后，靠个人经验写代码无法稳定控制质量、进度、成本和维护。

\subsection{1.3 软件工程的定义}

IEEE 定义的关键词：

\begin{itemize}
\item 系统化
\item 规范化
\item 可度量
\item 开发、运行和维护
\end{itemize}

可写为：

\begin{quote}
软件工程是把系统化、规范化、可度量的方法应用于软件的开发、运行和维护，并研究这些方法的学科。
\end{quote}

软件工程解决的不是“怎么写出一段代码”，而是“如何在多人协作、长期演化、质量受控的条件下交付软件系统”。

\subsection{1.4 软件危机与 Ariane 5 案例}

软件危机的标志性节点是 1968 年 NATO 软件工程会议。

常见问题：

\begin{itemize}
\item 项目延期
\item 成本超支
\item 软件质量低
\item 难以维护
\item 文档缺失
\item 需求不断变化
\item 开发过程不可控
\end{itemize}

Ariane 5 案例的答题重点：

\begin{itemize}
\item 代码本身来自已有系统，不等于在新环境中正确。
\item 隐含的数值范围假设没有被充分记录和验证。
\item 复用没有配套完成需求、环境、接口和异常条件验证。
\item 失败属于需求假设、集成验证、配置和变更管理问题，不只是编码问题。
\end{itemize}

\subsection{1.5 SWEBOK 知识域}

传统知识域覆盖需求、设计、构造、测试、维护、配置管理、工程管理、过程、模型方法、质量、职业实践、经济学、计算基础、数学基础、工程基础等。

SWEBOK V4 增强关注：

\begin{itemize}
\item 软件体系结构
\item 软件工程运维
\item 软件安全
\item AI for SE：用 AI 辅助软件工程
\item SE for AI：用软件工程方法建设 AI 系统
\end{itemize}

答题时可以用一句话收束：

\begin{quote}
AI 增强了需求整理、设计草图、代码生成和审查效率，但工程纪律仍由需求验证、接口约束、测试、配置管理和人工评审保证。
\end{quote}

\bigskip\hrule\bigskip

\subsection{1.6 项目与项目管理}

PMBOK 中项目的核心特征：

\begin{itemize}
\item 临时性：有明确开始和结束。
\item 独特性：交付独特产品、服务或成果。
\end{itemize}

软件项目管理要控制三个基本约束：

\begin{itemize}
\item 范围
\item 时间
\item 成本
\end{itemize}

三者共同影响质量。范围扩大、时间压缩或成本不足，都会传导到质量风险。

\subsection{1.7 项目管理五大过程组}

按课程表达：

\begin{enumerate}
\item 项目启动
\item 项目计划
\item 项目执行
\item 跟踪与控制
\item 项目收尾
\end{enumerate}

常见管理活动：

\begin{itemize}
\item 项目计划
\item 团队管理
\item 成本管理
\item 软件质量保证
\item 度量
\item 过程管理
\item 进度管理
\item 风险管理
\item 配置管理
\end{itemize}

\subsection{1.8 风险管理}

风险管理循环：

\begin{enumerate}
\item 风险识别
\item 风险分析
\item 优先级排序
\item 风险应对
\item 风险监控
\end{enumerate}

常见风险类型：

\begin{itemize}
\item 技术风险：新框架、新平台、性能瓶颈、兼容性问题。
\item 人员风险：关键成员离开、经验不足、沟通失败。
\item 需求风险：需求不清、范围膨胀、利益相关者冲突。
\item 进度风险：估算失真、依赖延期、任务拆分不合理。
\item 外部风险：政策、供应商、硬件、网络、第三方接口变化。
\end{itemize}

答题模板：

\begin{quote}
先列风险，再写影响，再写应对措施。  
例如：第三方支付接口变更属于外部和技术风险，会影响充值功能上线；应对措施是提前建立接口契约、Mock 服务、回归测试和变更通知机制。
\end{quote}

\subsection{1.9 软件质量保证 SQA}

SQA 关注“过程和产品是否满足质量要求”。

常见活动：

\begin{itemize}
\item 代码审查
\item 静态分析
\item 测试
\item 过程审计
\item 质量度量
\item 阶段评审
\end{itemize}

里程碑质量活动：

\begin{itemize}
\item 需求阶段：需求评审、需求度量。
\item 体系结构阶段：架构评审、集成测试策略。
\item 详细设计阶段：设计评审、设计度量、集成测试准备。
\item 构造阶段：代码审查、代码度量、单元测试。
\item 测试阶段：缺陷度量、覆盖率分析、回归测试。
\end{itemize}

\subsection{1.10 软件配置管理 SCM}

SCM 的目标：

\begin{itemize}
\item 识别配置项
\item 控制变更
\item 记录和报告配置状态
\item 审计配置
\item 保证产品与规格一致
\end{itemize}

核心术语：

\begin{itemize}
\item 配置项 CI：纳入配置管理的项目制品，如需求文档、设计文档、源代码、测试脚本、部署脚本、数据库脚本。
\item 基线 Baseline：经过正式评审和批准的规格或产品，后续变更必须走正式变更控制。
\end{itemize}

SCM 常见活动：

\begin{enumerate}
\item 配置项识别
\item 版本管理
\item 变更控制
\item 配置审计
\item 状态报告
\item 发布管理
\end{enumerate}

\subsection{1.11 变更控制流程}

标准作答流程：

\begin{enumerate}
\item 提交变更请求 CR。
\item 进行影响分析：影响需求、设计、代码、测试、进度、成本和风险。
\item 进行风险评估。
\item CCB 审查并批准或拒绝。
\item 制定回滚计划。
\item 执行变更。
\item 更新基线。
\item 上线验证。
\item 更新状态报告。
\end{enumerate}

Knight Capital 案例的答题重点：

\begin{itemize}
\item 部署遗漏一台服务器。
\item 新旧代码路径并存导致异常交易。
\item 这是发布管理、配置管理和上线验证失败。
\item 技术错误背后是过程控制缺失。
\end{itemize}

\bigskip\hrule\bigskip

\bigskip\hrule\bigskip

\section{对应考试模板}

\subsection{7.12 配置管理题模板}

题目出现上线后版本不一致、漏部署、旧代码被触发：

\small
\begin{verbatim}
问题属于配置管理和发布管理失败。
原因：配置项识别不完整，版本状态未核对，发布前未做配置审计和上线验证。
改进：
1. 将源码、脚本、配置、依赖和部署清单纳入配置项。
2. 建立发布基线。
3. 变更走 CR、影响分析、CCB 审批、执行、审计、状态报告。
4. 发布前核对所有服务器版本。
5. 发布后执行冒烟测试和回滚预案。
\end{verbatim}
\normalsize

\bigskip\hrule\bigskip

\bigskip\hrule\bigskip

\section{快速背诵清单}

\subsection{软件工程}

\begin{itemize}
\item 软件 = 程序 + 数据 + 文档 + 知识。
\item 软件工程 = 系统化、规范化、可度量的方法。
\item 软件危机来自规模、复杂度、变化和协作失控。
\item AI 加速工程活动，但不能替代验证、评审和配置管理。
\end{itemize}

\subsection{项目管理}

\begin{itemize}
\item 项目：临时性、独特性。
\item 三约束：范围、时间、成本，影响质量。
\item 五过程组：启动、计划、执行、跟踪控制、收尾。
\item 风险管理：识别、分析、排序、应对、监控。
\item SCM：配置项、基线、版本、变更、审计、状态报告、发布。
\end{itemize}


\end{document}
